\chapter{Discussion}
\label{ch:discussion}

This chapter interprets the results of Chapter~\ref{ch:evaluation}.
Section~\ref{sec:discussion_core} states the central finding and
Section~\ref{sec:discussion_h1h6} reads H1 and H6 together.
Sections~\ref{sec:external_interpretability} to~\ref{sec:ablation_establishes}
set out the scope of the external estimate, the anchor-independent label and
the feature ablation. Section~\ref{sec:implications} draws out implications for
how studies of this kind are reported, and Section~\ref{sec:limitations} sets
out the limitations.


\section{The Core Finding}
\label{sec:discussion_core}

The central finding concerns outcome construction. Under the pre-registered
conjunction-based timing rule, no Sepsis-3 onset can occur before the treatment
anchor (Proposition~\ref{prop:h3}). Applying the standard
$\min(t_{\text{susp}}, t_{\text{SOFA}})$ timing to the same cohort, whose
septic membership is cohort-wide identical, recovered \pcPreTreatOnsetsVarF{}
pre-treatment onsets, separable from their surrounding at-risk hours at
\pcPreTreatAurocPrimaryF{} AUROC on the onset-hour estimand. Whether those
onsets are predictable in advance cannot be assessed here, because predictors
are binned to the onset hour (Limitation~\ref{lim:current_hour}).

This finding has an analytic and a measured part. Analytically, any study that
uses a conjunction-based onset rule and reports pre-treatment performance is
reporting on an empty set, and a study that does not restrict the window is,
to that extent, scoring discrimination around a clinician-defined decision.
Empirically, the operationalisation rather than the data removes the
pre-treatment onsets. Neither part supports the proposition that no patient
develops sepsis before clinicians act, which external evidence contradicts:
the Epic Sepsis Model v2 evaluation \cite{wong2026epic} reports median lead
times of 1.4--7.1 hours ahead of clinician recognition across four sites.
Sepsis-3 \emph{membership}, however, depends on clinical action under every
timing rule, since without a qualifying antibiotic--culture pair there is no
label, so Kamran et al.'s \cite{kamran2024evaluation} concern applies to which
patients can enter the positive class.

The Utility Score stays close to the no-alert strategy for every model. This
reflects both the low hourly event rate and the scoring rule's cumulative
penalty for false-alert hours, which is charged for every alerting hour while
only the first timely alert is credited; the study does not separate their
individual contributions, nor test whether the temporal split contributes
through era-related differences in recording. AUROC remains in the high 0.7s
under the same conditions because the ranking of person-hours is preserved
even when the achievable alert volume is not clinically usable. Wang et al.\
\cite{wang2025srpm} documented this divergence across 91 models moving from
internal to external validation (Section~\ref{sec:auroc_illusion}); here it is
already present internally, under a time-based split, where utility is close
to the no-alert baseline before any external validation.


\section{H1 and H6 Together}
\label{sec:discussion_h1h6}

H1 and H6 address different levels of comparison. H1 compares the label spread
with the spread across six complete model specifications, which captures the
large gap between flexible fitted models and simpler comparators: the
specification spread \pcSpreadModelClass{} exceeds the label spread
\pcSpreadLabel{}, and the interval on the difference excludes zero. H6
compares two flexible models fitted on the same predictors and shows a much
smaller difference, \pcHSixAbsEst{} AUROC in the primary model's favour, with
GBT not better by more than the pre-registered margin. In this study,
therefore, the difference between the two flexible estimators was small
relative to the difference between fitted models and simpler comparators,
while the label changed which patients were counted far more than it changed
discrimination.

The H6 result also bears on the statistics-versus-machine-learning question of
Section~\ref{sec:false_dichotomy}. Christodoulou et al.'s
\cite{christodoulou2019cox} equivalence finding concerns binary outcomes on
tabular data, and Fagerstr{\"o}m et al.\ \cite{fagerstrom2019lisep} reported a
deep model detecting septic-shock onset up to 20 hours earlier than a Cox
baseline. Here, with the label, the split, the features and the evaluation
held fixed, the primary discrete-time hazard model is not outperformed by
GBT. This does not contradict Fagerstr{\"o}m et al.: their comparator was a
static Cox specification, the same family as the Cox benchmark retained here,
which is also well behind the fitted models in this study (\pcAurocCoxB{}
against \pcAurocPrimaryB).

The intervals add one thing the point estimates could not. The label spread's
interval [\pcSpreadLabelCiLo, \pcSpreadLabelCiHi] and the specification
spread's [\pcSpreadModelClassCiLo, \pcSpreadModelClassCiHi] do not overlap, so
H1's reversal is not attributable to resampling noise. The label effect is
real, since its interval excludes zero, but its main consequence is on the
cohort, the event count and every denominator computed on it, which AUROC does
not capture.


\section{Interpreting the External Validation}
\label{sec:external_interpretability}

The interpretability floor adopted here is \pcExtMinEvents{} events, following
the external-validation sample-size guidance of Riley et al.\
\cite{riley2019sample} and Collins et al.\ \cite{collins2024tripodai}: below it
the confidence interval on an external AUROC is wider than the degradation the
study is trying to detect. All three Sepsis-3 variants fall below it, with
\pcExtNEventsB{} events under the primary label, so the external Sepsis-3
degradations are stated as directions rather than magnitudes, and the H5
interval is the formal expression of that limit. External AUPRC and Utility
Score are not interpreted at this event count.

The microbiology-covered sub-cohort is also selected by documentation: a
culture is ordered when infection is already suspected, so the covered stays
are enriched for patients a clinician had begun to investigate, and the
external estimate describes transport into that sub-population rather than
into eICU-CRD as a whole (Limitation~\ref{lim:label_transport}).

The antibiotic-only arm clears the floor with two orders of magnitude more
events, but it labels a broader target, treated suspicion of infection without
the culture limb, and is never pooled with the Sepsis-3 arm. Its degradation
has the same sign and label ordering as the Sepsis-3 arm's and a somewhat
larger magnitude. Because the two arms label different targets, that agreement
is descriptive. Accordingly, the external analysis supports a direction of
effect but not a precise estimate for Sepsis-3. The same arm shows that the
internal ordering of comparators does not transport: GBT scores
\pcExtAbxAurocGbtB{} under Variant B, close to chance under the two narrower
labels, while NEWS2 scores \pcExtAbxAurocNewsB, and a bedside score ranks ahead
of GBT under every label in this arm. The two flexible estimators are close
internally, but only the primary model stays ahead of NEWS2 externally. The
arm's event count supports this ordering but not the size of any gap.

The Variant D arm is the one external estimate above the floor for its own
target (Section~\ref{sec:variant_d_establishes}); it provides precise
transport estimates for deterioration but does not improve inference for
Sepsis-3. The Sepsis-3 point estimate, \pcHFiveEst, is about three-quarters of
the 0.085 penalty Moor et al.\ \cite{moor2023review} report across four
databases. Unlike Moor et al., this study cannot attribute the loss, because
their harmonised databases shared labelling inputs and eICU-CRD does not.


\section{Scope of the Variant D Results}
\label{sec:variant_d_establishes}

Variant D's label is a threshold function of recorded physiology, and five of
the six SOFA components are model inputs, so its discrimination is partly
self-fulfilling (Section~\ref{sec:alt_labels}). Its AUROC \emph{level}, internal
or external, is therefore not comparable with Variant B's and is not evidence
about sepsis prediction. This caveat applies throughout this section. Within
it, Variant D supports four conclusions.

\begin{enumerate}
    \item \textbf{The label sensitivity is not wholly definitional.} At
    $\kappa = \pcKappaVsBD$ and a \pcPctExactMatchD\% onset match, an
    anchor-independent label disagrees with Sepsis-3 about both membership and
    timing, so the H1 label result is not only an artefact of how one
    definition is read.

    \item \textbf{The alert burden remains high under every label; the
    utility ceiling does not.} The best achievable Utility Score rises from
    \pcUtilityMaxGbtB{} under Variant B to \pcUtilityMaxGbtE{} under E and
    \pcUtilityMaxGbtD{} under D, with the primary model following the same
    pattern. Even D's best configuration carries \pcUtilityOptBurdenGbtD{}
    false alerts per true alert, several times the 1.4 benchmark of Moor et
    al.\ \cite{moor2023review}. Contribution C3 therefore holds across labels
    in its alarm-volume form but not as a utility ceiling near zero; the rise
    is itself expected for a label built from recorded physiology.

    \item \textbf{The two metrics order the models differently, and the
    ordering is stable across labels.} At its own best threshold GBT leads on
    the Utility Score under B, D and E alike, while on AUROC the primary model
    leads under B and E and GBT only under D. Changing the label does not
    reverse either ordering; changing the metric does. This illustrates the
    concern behind H5, that the verdict depends on the metric reported, without
    testing H5's internal-to-external contrast.

    \item \textbf{The external event-rate gap is largely attributable to the
    anchor.} The Sepsis-3 arm yields an external event rate
    \pcExtRateRatioB{} times the development cohort's. Applied to all
    \pcExtAltNStays{} eICU-CRD stays (Protocol Amendment~9), Variant D gives a
    ratio of \pcExtAltRateRatio. Because the two labels are applied to the
    same two cohorts, a population difference large enough to produce the
    anchored ratio would also have produced a comparable shortfall under the
    anchor-independent label, and it does not. The reduction therefore
    suggests that label transport accounts for a substantial part of the
    original discrepancy. It bounds the anchor's contribution rather than
    isolating it: roughly a quarter of the development cohort's rate remains
    unaccounted for, and case mix, charting frequency, the five-component
    external SOFA rise and differences in illness severity all sit inside that
    residual.
\end{enumerate}

The large Variant D event count also supports the transport \emph{drop} with
an interval (Protocol Amendment 10): both models lose AUROC on transport, and
both intervals exclude zero (Table~\ref{tab:external_altlabel}). This is the
study's one demonstrated transport loss. It concerns a deterioration rule, not
Sepsis-3, and cannot be pooled with the anchored arm. It does show that the two
databases are not similar enough for a frozen model to move between them at no
cost, which is consistent with the Sepsis-3 arm being unable to resolve a
degradation of the size the literature anticipates on \pcExtNEventsB{} events.


\section{Interpreting the Ablation}
\label{sec:ablation_establishes}

The ablation is exploratory and run on one cohort, and with \pcHTwoNTerms{}
covariates the stability counts are small enough that the comparison is
descriptive. Its interpretation has three parts.

First, the explicit action-derived features contribute little to AUROC.
Removing them changes discrimination by \pcAblDeltaAurocPrimaryB{} for the
primary model under the primary label and by less than 0.01 in all six
comparisons. \pcAblNExcludeZero{} of the \pcAblNContrasts{} paired intervals
exclude zero, and no margin of practical equivalence was pre-registered, so
this is a statement about magnitude, not absence; a feature that moves AUROC
little may still be the reason a particular alert fires.

Second, these features account for most of the cross-label coefficient
instability: \pcHTwoPctUnstable\% of terms are flagged in the full fit against
\pcAblHTwoPctUnstable\% in the ablated one. A published effect estimate for a
treatment or order indicator from a model of this kind therefore reflects the
labelling regime as much as the patient.

Third, removing them does not remove all clinician-action information. The
retained laboratory values are forward-filled from tests somebody chose to
order, and vital signs are recorded on schedules that reflect how closely a
patient is watched, so the ablation measures the effect of removing the
explicit traces only (Limitation~\ref{lim:features}).


\section{Implications for How These Studies Are Reported}
\label{sec:implications}

The findings bear less on which model to build than on what a sepsis prediction
study needs to report before its headline number can be interpreted.

\paragraph{Report label agreement, not only performance sensitivity.}
Performance across label variants should not be reported as the only
label-sensitivity analysis. Here, label choice moved AUROC by \pcSpreadLabel{}
while the two narrowest readings agreed at $\kappa = \pcKappaVsBA$ despite
near-identical prevalence, so every count, rate and denominator changed
substantially while discrimination barely moved. Agreement between the
variants is where the variation lies, and AUROC does not show it.

\paragraph{State the onset rule and its consequence for the evaluation
window.} Two studies can share a cohort, a definition and a model and still
differ on whether pre-onset prediction is measurable, purely through the choice
between $t_{\text{susp}}$ and $\min(t_{\text{susp}}, t_{\text{SOFA}})$. That
choice is often left implicit in a reference to Sepsis-3. It should be stated
in the methods with its consequence for the evaluation window, because under a
conjunction rule an empty pre-treatment window is not evidence of anything
(Proposition~\ref{prop:h3}).

\paragraph{Report discrimination and operating-point burden together.} At the
event rate here the conventional 0.5 cut-off separates models by calibration
rather than clinical value. A threshold-dependent metric is interpretable only
as a swept ceiling with the attaining operating point reported beside it,
including its alert burden (\pcUtilityOptBurdenGbtB{} false alerts per true
alert at the best configuration under the primary label, against the 1.4
benchmark of Moor et al.\ \cite{moor2023review}). Alert burden is rarely
reported among the 91 models Wang et al.\ \cite{wang2025srpm} reviewed.

\paragraph{Audit label-input coverage before external validation.} In this
study the main constraint on the external estimate was whether the label's
inputs existed in the destination database, which neither a reported AUROC
nor a reported cohort size reveals. A study transporting a treatment-anchored
label should report the coverage of each labelling input in the destination
cohort, and the event count an estimate rests on, before reporting performance.

\paragraph{Interpret comparisons between flexible estimators in context.}
Within this study, comparisons between flexible estimators were less
informative than comparisons of label membership, onset timing and alert
burden: the two flexible estimators differ by \pcHSixAbsEst{} AUROC against
\pcSpreadModelClass{} across the six specifications. Two estimators are a
narrow base for a general claim, since no deep temporal model was fitted and
GBT was not tuned (Limitation~\ref{lim:comparators}).

\paragraph{Declare subgroup event floors in advance.} At a low event rate most
demographic strata carry too few events for a subgroup AUROC to be
interpretable, and an interval reaching the chance line is easily read as a
null result rather than an absent measurement. Declaring the floor before
estimation, tabulating below-floor strata rather than suppressing them, and
excluding them from the corrected family keeps that distinction visible.


\section{Limitations}
\label{sec:limitations}

The limitations below are not of equal weight, and three bound the central
claims most directly: the discrimination figures measure detection at the
recorded onset hour rather than advance warning
(Limitation~\ref{lim:current_hour}); part of the label-variant spread is
definitional, because the pre-registered variants share one treatment-based
anchor (Limitation~\ref{lim:label_construction}); and the strongest positive
result, H6, is non-inferiority against a single untuned comparator whose
boosting-round selection was corrected under Protocol Amendment~8
(Limitation~\ref{lim:comparators}). The remainder qualify the scope of
particular analyses rather than the direction of the main findings.

\begin{enumerate}
    \item \textbf{Development cohort and population
    selection.}\label{lim:cohort_selection} All models were developed on
    MIMIC-IV, from a single academic medical centre, and the findings apply to
    the Sepsis-3 definition on that database; clinical practice, documentation
    and antibiotic prescribing elsewhere may differ, and SEP-1 and the CDC
    Adult Sepsis Event were not implemented. The base cohort is also selected:
    only the first ICU stay of each patient and stays of at least four hours
    are retained (Section~\ref{sec:cohort}). Repeat admissions differ
    systematically from first admissions, and short stays mix rapid deaths,
    rapid transfers and low-acuity observation admissions; the two
    restrictions pull acuity in opposite directions, so their net effect is not
    signed here. The population is the subset of adult ICU admissions for
    which an hourly real-time task is well posed.

    \item \textbf{External label transport and external sample
    size.}\label{lim:label_transport} The Sepsis-3 anchor requires a culture,
    which \texttt{microLab} documents for only \pcExtMicroCoverage\% of the
    \pcExtNCohortStays{} eICU-CRD stays, so the primary external analysis rests
    on \pcExtNEventsB{} events under the primary label, below the
    \pcExtMinEvents{}-event floor. More data of the same kind would not remove
    this; it would require a destination database whose microbiology recording
    matches MIMIC-IV's. The covered sub-cohort is documentation-selected, with
    a direction of selection that cannot be signed, and the external
    degradation is not decomposed into case mix, measurement frequency,
    presentation or residual label imperfection. NEWS2 is also not scored
    identically in the two cohorts: in MIMIC-IV a person-hour with no recorded
    Glasgow Coma Scale after forward-fill is charged three points as ``not
    alert'', whereas in eICU-CRD, which lacks the variable, it is charged none,
    so NEWS2 thresholds, specificity and alert burden should be compared within
    a cohort, not across the pair. No decomposition, verdict or primary-model
    result depends on NEWS2.

    \item \label{lim:current_hour}\textbf{Discrimination is measured on the
    current-hour hazard, not on a six-hour horizon.} Every AUROC, AUPRC and
    calibration figure scores the current-hour hazard against an indicator
    marking the onset hour among at-risk hours (Section~\ref{sec:framing}).
    The figures are therefore not directly comparable with published AUROCs for
    binary ``onset within $k$ hours'' labels, whose positives are more numerous
    and whose task is easier to score well on, and comparisons with such
    figures in Chapter~\ref{ch:review} are indicative only. Covariates are
    binned to the hour, so a measurement charted after the onset timestamp
    within the onset hour still enters that hour's prediction: the figures
    measure detection at the recorded onset hour, not anticipation, and the
    lead-time distribution, which speaks to anticipation, is uninformative at
    these alert rates. A design that lagged every predictor by one hour, or
    applied a sub-hour timestamp cut-off, would separate the two and is not
    implemented here. A cumulative six-hour risk, which a deployed instrument
    would display, is also not reported.

    \item \textbf{Label construction.}\label{lim:label_construction} Part of the label-variant spread is
    definitional: Variants A--C share the culture-plus-antibiotic anchor, so
    moving the antibiotic--culture windows must move the onsets. The
    pre-registered onset rule is also not the standard one, so the
    pre-registered family cannot address onset timing on its own, H3 among
    it. The alternative labels that address these points are post-hoc
    (Protocol Amendments 2 and 4) and rest on few variants. Variant F uses the
    full SOFA score, so its onsets are free of the antibiotic--culture anchor
    but not of treatment in general, and its $t_{\text{SOFA}}$ is searched only
    inside the 72-hour panel. Variant D is an operational deterioration label
    built for this comparison, not a validated construct, its onsets have no
    external adjudication, and its overlap with the model inputs makes its
    discrimination level uninterpretable.

    \item \textbf{Feature construction.}\label{lim:features} The feature set
    cannot be fully purged of clinician action: the ablation removes the
    explicit action-derived predictors, but retained laboratory values and
    vital signs still reflect ordering and monitoring decisions. Clinical
    notes, procedure codes and ventilator settings were not included. The
    covariate panel is also single-analyst work with no independent
    re-implementation; declared-but-absent features and implausible recorded
    values are guarded against, by an error on any missing feature and a
    plausibility range on all nineteen physiological variables before
    forward-fill, but neither guard substitutes for independent replication.

    \item \textbf{Comparator coverage and the scope of
    H6.}\label{lim:comparators} A deep survival comparator was planned but not
    built, so H6 establishes non-inferiority against GBT only, not against
    every flexible estimator. GBT was not tuned, which keeps the search budget
    equal across the H6 comparison but may place it below its attainable
    performance. The contrast supports one-sided non-inferiority, not two-sided
    equivalence (Section~\ref{sec:h6}). GBT's boosting-round count is chosen on
    a holdout from the training stays (Protocol Amendment 8); under the earlier
    procedure it was chosen on the test set, which favoured the comparator, and
    the correction moved H6 in the primary model's favour, so H6 is a finding
    the uncorrected fits already supported and the correction strengthened.

    \item \textbf{Penalised estimation and H2 uncertainty.} The unpenalised
    likelihood is completely separated, so the reported fits are
    ridge-penalised; test AUROC varies by less than 0.02 across a 20-point
    penalty path, but the coefficients are shrunk estimates and the H2
    counts are properties of the penalised solution. The cluster-robust
    covariance is that of the penalised estimator, so Wald-type
    interpretation is approximate, and it bears on coefficients only: H2's
    count has no sampling distribution. The Amendment 7 diagnostic shows that
    \pcHTwoNNoise{} of the \pcHTwoNFlagged{} flagged covariates move by less
    than their own estimation noise, so the count is a loose upper bound, but
    the diagnostic is not a test. Performance intervals come from the cluster
    bootstrap and are unaffected.

    \item \textbf{Threshold selection and alarm implementation.} Operating
    points and maximum Utility Scores are selected on the test set on which
    they are evaluated, so they are optimistic upper bounds; selecting
    thresholds on the training years would give a lower deployment estimate.
    External thresholds are re-derived on eICU-CRD, which hides the further
    loss a transferred threshold would incur. False-alarm episodes are defined
    by contiguity rather than acknowledgement; suppression and escalation logic
    would reduce them, so the alert-hour and episode rates bound the true
    burden from above and below.

    \item \textbf{Subgroup analysis.}\label{lim:subgroup} Only
    \pcNSubgroupsAboveFloor{} of the racial strata clear the
    \pcMinSubgroupEvents{}-event floor and no contrast in family E1 survives
    correction, so the data are consistent both with no disparity and with
    material disparities. The strata are also administrative: the
    \texttt{race} field mixes broad and granular codes from the same
    population, so the reference level is a residual category, and two levels
    record the absence of an answer. The label-sensitivity result is better
    powered but inherits the construct limitation, and its two significant
    racial strata are the missingness categories. A documentation-based
    explanation, in which an incomplete registration record accompanies an
    incomplete clinical record and so makes a label built on two documented
    clinician actions less stable, is assumed rather than measured; comparing
    charting density or laboratory coverage between strata would test it.

    \item \textbf{Literature-search scope.} The search in
    Section~\ref{sec:search_strategy} is structured and verified but not
    systematic: no title-and-abstract screening of a full deduplicated yield
    was performed and no second reviewer screened independently. The novelty
    claims for contributions C1 and C5 concern combinations of analyses, each
    with precedent, and a systematic review could surface a prior study
    combining them.

    \item \textbf{Post-hoc amendments and hypothesis evaluability.} Ten
    protocol amendments were introduced after the initial analyses had been
    inspected. They are dated and disclosed in the pre-registration, and their
    analyses are tagged exploratory or descriptive (Appendix~\ref{sec:register}),
    but the confirmatory interpretation should be read as a pre-specified
    hypothesis set analysed under a retrospectively specified inference plan
    (Amendment 1). Amendment 8 corrected test-set use in GBT early stopping and
    changed the GBT estimates; no amendment altered an H1--H6 estimand. Three
    others warrant a note. Amendment 6 reduced family E1 from thirty computed
    contrasts to six, a departure from the rule that the denominator is the
    number of contrasts computed; the outcome-blind event floor removed the one
    contrast that had cleared correction, resting on twelve events, so the
    change runs against this study, and none of the six retained contrasts
    clears correction under either denominator. Amendment 7's diagnostic
    favours this thesis's argument and is computed from the same fits as the
    ablation, so it is not independent evidence. Amendment 10 attached
    intervals to three quantities outside the corrected families:
    \pcAblNExcludeZero{} of the \pcAblNContrasts{} ablation contrasts and one
    cross-label calibration-intercept contrast exclude zero. Of the
    hypotheses, H2 has no test statistic, H3 is degenerate by construction,
    and H5 is tested on a substituted estimand; that substitution is a
    departure from the protocol, and H5's verdict is weaker evidence than H1's
    or H4's for that reason.
\end{enumerate}
